\documentclass{article}
\usepackage{hyperref}
\usepackage{nopageno}
\usepackage{graphicx}

\title{\textbf{It Doesn't Work!}}
\date{}
\author{}

\begin{document}
\maketitle
\vspace{-50pt}

% amphibrachic tetrameter
\begin{center}
The code has stopped working with no explanation.\\
It doesn't make sense and you cry in frustration.\\
Just take a deep breath and use these three suggestions,\\
so you will be able to ask the right questions.\\
\end{center}

\section*{Ask yourself:}
\begin{itemize}
\item \textbf{\emph{Why} do you need it to work?}\\
What is the context?
What is the goal?

\item \textbf{How do you \emph{show} that it doesn't work?}\\
What are the steps to reproduce your problem?
How will you know that the problem is resolved?

\item \textbf{What \emph{does} work?}\\
How far are you from a working state?
What similar cases work?
\end{itemize}

Then, \underline{ask for help, without fear!}

\vspace{10pt}

\section*{Helpful Practices}
A good workflow will help to avoid problems in the first place.
\begin{itemize}
\item Work on one problem at a time.
\item Write small, simple, and quick-to-run test cases.
\item Test often.
\item Visualize often.
\item Save working states often.
\item Plan for reproducibility.
\end{itemize}
Version control (e.g. \texttt{git}) is often useful.

\vspace{10pt}
\begin{center}
{\small
By \href{https://www.patricksanan.com}{Patrick Sanan}, 2019.
This document is released under a Creative Commons license and is available at
\href{https://github.com/psanan/it_doesnt_work}{\texttt{github.com/psanan/it\_doesnt\_work}}.
Version 1.0 .\\
\vspace{10pt}
\includegraphics[width=50px]{by-nc-sa.png}
}
\end{center}

\end{document}
